Stochastic MPC extends deterministic MPC to handle uncertainties
modeled as random variables, such as additive noise or parameter
variations. A core feature is replacing hard constraints with chance
constraints to account for uncertainty.
\[
  \mathbb{P}(x_k \in \mathcal{X}) \geq 1 - \epsilon, \quad
  \mathbb{P}(u_k \in \mathcal{U}) \geq 1 - \delta,
\]
This softens constraints, improving feasibility over robust MPC.
Chance constraints can be joint or separate (per state/input or timestep).

\textbf{Mean Dynamics}
$\bar{x}_{k+1} = A\bar{x}_k + B\bar{u}_k$ (Deterministic)

\textbf{Variance Dynamics}
$\Sigma_{k+1} = (A+BK)\Sigma_k(A+BK)^\top + \Sigma_w$

The variance $\Sigma_k$ can be precomputed offline
as it does not depend on the actual state measurements.

For general distributions or polytopic sets, we use
\textbf{Probabilistic Reachable Set}
$\mathcal{F}_k^p$ such that
$\text{Pr}(e(k) \in \mathcal{F}_k^p) \geq p$,
where $e_k = x_k - \bar{x}_k$.

The constraint is satisfied if:
$\bar{x}_k \in \mathcal{X} \ominus \mathcal{F}_k^p$

For Gaussian noise ($w_i \sim \mathcal{N}(0, \Sigma_w)$),
$\mathcal{F}_k^p$ can be an ellipsoid: $\{e \mid e^\top \Sigma_k^{-1}
e \leq \chi_{n_x}^{-2}(p)\}$, where $\chi_{n_x}^{-2}$ is the inverse
chi-squared quantile.

\begin{sstTitleBox}{
    Stochastic MPC
  }
  \begin{sstOnlyFrame}
    \[J = \mathbb{E} \left[ \|x_N\|_P^2 + \sum_{k=0}^{N-1}
    \|x_k\|_Q^2 + \|u_k\|_R^2 \right]\]
    Using the property $\mathbb{E}[x^\top Q x] = \bar{x}^\top Q
    \bar{x} + \text{tr}(Q \Sigma_x)$, the cost decomposes into:

    \textbf{Nominal Cost}
    A standard quadratic cost on the mean
    trajectories $(\bar{x}, \bar{u})$.

    \textbf{Constant Offset}
    A term involving $\sum \text{tr}((Q + K^\top R K) \Sigma_k) +
    \text{tr}(P \Sigma_N)$,
    which depends only on the noise statistics and the gain $K$.

    \textbf{Conclusion}
    For a fixed $K$, optimizing the expected cost is equivalent
    to optimizing the nominal cost of the mean system.
  \end{sstOnlyFrame}
  \begin{sstOnlyFrame}
    \textbf{System dynamics}
    $x_{i+1} = A x_i + B u_i + w_i$ (probabilistic)\\

    \textbf{Chance constraints}
    $\mathbb{P}(x_i \in \mathcal{X}) \geq 1 - \epsilon$,
    $\mathbb{P}(u_i \in \mathcal{U}) \geq 1 - \delta$\\

    \textbf{Initial condition}
    $x_0 = x_k$

    \textbf{Terminal set}
    $\mathbb{P}(x_N \in \mathcal{X}_f) \geq 1 - \epsilon$ (probabilistically)
  \end{sstOnlyFrame}
  \begin{sstOnlyFrame}
    \textbf{Tube-like Feedback Policy}
    (to control spread of trajectories)
    \[u(k) = \bar{u}_k + K(x(k) - \bar{x}_k)\]
    For tractability, approximations like affine feedback policies $u_i =
    K x_i + v_i$ (where $v_i$ is optimized) are used to close the loop in
    predictions. This contrasts with open-loop predictions, which
    ignore feedback in the horizon and may lead to conservative results.
  \end{sstOnlyFrame}
\end{sstTitleBox}

\subsubsection{Feasibility}
Large (unbounded) disturbances can drive the system into regions
where the optimization problem becomes infeasible at the next time step.

\subsubsection{Asymptotic Average Performance Bound}
For a stabilizing gain $K$ satisfying the Lyapunov equation $P =
(A+BK)^\top P (A+BK) + (Q + K^\top R K)$,
$\lim_{T \to \infty} \frac{1}{T} \sum \mathbb{E}[l(x,u)] \leq
\text{tr}(P\Sigma_w)$.

\subsubsection{Scenario Approach}
For non-Gaussian uncertainties, sample $M$ disturbance scenarios
$w^{(j)}$ and enforce constraints for each. Choose $M$ based on
desired violation probability (e.g., via Chernoff bounds). This is
computationally intensive but distribution-free.

\textbf{Stochastic vs. Robust MPC}: Stochastic allows probabilistic
violations for better performance/feasibility; robust ensures
worst-case guarantees but can be conservative.

\subsection{Recovery Initialization}
If the measured state $x(k)$ is infeasible for the MPC, the solver is
initialized with the \textbf{predicted mean} from the previous step
($\bar{x}_{1|k-1}$) instead of the actual measurement. This ensures
recursive feasibility but sacrifices performance.

\subsection{Indirect Feedback SMPC}

A more advanced formulation where the nominal state $z(k)$
evolves independently of the measurement $x(k)$ in the constraints,
but $x(k)$ enters the \textbf{cost function}.

\textbf{- Pro}
Guarantees recursive feasibility and satisfaction
of closed-loop chance constraints.

\textbf{- Con}
The feedback is "weaker" than direct state-feedback MPC.

